\cleardoublepage
\phantomsection
\addcontentsline{toc}{section}{Введение}
\section*{Введение}

В статистической теории обучения~\cite{vapnik1998statistical,shalev2014understanding} обоснована принципиальная связь между сложностью семейства моделей, объемом обучающей выборки и способностью к обобщению: для надежного выбора модели требуется достаточное количество данных. Подавляющее большинство современных моделей машинного обучения являются параметрическими: они задаются конечным набором параметров, определяемых минимизацией эмпирической функции потерь на выборке. Поскольку функция потерь вычисляется по ограниченным данным, при изменении выборки меняется как множество ее минимизаторов, так и сама поверхность функции потерь в их окрестности. Это ставит вопрос о том, как геометрия оптимизационной поверхности стабилизируется при росте объема обучающей выборки и при каких условиях можно говорить о достаточности выборки для конкретной модели.

Эмпирические законы масштабирования~\cite{kaplan2020scaling,hoffmann2022chinchila} устанавливают степенные зависимости качества параметрических моделей от объема данных и числа параметров, однако они получены на основе масштабных экспериментов и в настоящее время не имеют единого теоретического обоснования. Параллельно ведутся исследования локальной геометрии ландшафта функции потерь~\cite{li2018visualizing,fort2019large,draxler2018essentially} и спектральных свойств матрицы Гессе в окрестности минимума~\cite{sagun2017hessian,ghorbani2019investigation,papyan2019measurements,papyan2020traces}, выявившие, что спектр Гессе глубоких моделей обладает выраженно неоднородной структурой с небольшим числом доминирующих собственных значений. Однако систематической связи между объемом обучающей выборки и стабилизацией самой поверхности функции потерь, согласованной с этой эмпирически наблюдаемой анизотропией, в литературе не предложено.

\textbf{Актуальность темы.} Существующие способы количественного измерения стабилизации ландшафта функции потерь "--- одноточечная оценка приращения функции потерь в фиксированной точке и изотропное среднеквадратичное усреднение по гауссовскому распределению точек в полном пространстве параметров "--- являются частными случаями более общей конструкции с функцией предпочтения точек в пространстве параметров. Изотропное среднеквадратичное усреднение приводит к верхней оценке порядка $\mathcal O(k^{-2})$, в которой, однако, доминирующая константа зависит от размерности пространства параметров $N$ как $N^2+2N$. В нейронных сетях современных размеров, где $N$ достигает $10^7$"--$10^9$, такая зависимость делает количественное применение существующей оценки невозможным.

В то же время эмпирически наблюдаемая анизотропия спектра матрицы Гессе указывает на то, что существенное изменение поверхности функции потерь при малых возмущениях сосредоточено в подпространстве сравнительно небольшой размерности, натянутом на ведущие собственные векторы матрицы Гессе. Это мотивирует постановку, в которой пробная мера согласована с геометрией модели: переход от изотропного усреднения в полном пространстве к усреднению в подпространстве наибольшей кривизны. Разработка соответствующих критериев, доказательство их сходимости и построение масштабируемых алгоритмов оценивания являются актуальной задачей, поскольку именно они позволяют получить количественные критерии достаточности обучающей выборки, применимые к моделям современных размеров.

\textbf{Объект исследования} "--- параметрические модели машинного обучения и оптимизационные задачи минимизации эмпирической функции потерь, возникающие при их обучении на конечных выборках.

\textbf{Предмет исследования} "--- локальная геометрия поверхности эмпирической функции потерь в окрестности ее минимума и ее стабилизация при увеличении объема обучающей выборки.

\textbf{Цель работы} "--- получить количественные оценки достаточного размера обучающей выборки параметрических моделей машинного обучения на основе анализа стабилизации поверхности эмпирической функции потерь в окрестности ее минимума.

\textbf{Задачи работы.} Для достижения поставленной цели решаются следующие задачи:

\begin{enumerate}
    \item Провести анализ существующих подходов к измерению стабилизации ландшафта функции потерь в нейронных сетях и выявить ограничения постановок, основанных на одноточечной оценке и изотропном усреднении в полном пространстве параметров.
    \item Сформулировать общий критерий стабилизации поверхности функции потерь с функцией предпочтения точек в пространстве параметров, охватывающий существующие частные постановки в единой схеме.
    \item Получить верхние оценки скорости убывания критериев стабилизации при увеличении объема обучающей выборки, в том числе для критерия, ограниченного на подпространство наибольшей кривизны матрицы Гессе, и построить масштабируемые алгоритмы их вычисления.
    \item Провести вычислительную проверку полученных оценок на полносвязной нейронной сети и трансформерной языковой модели.
\end{enumerate}

\textbf{Методы исследования.} Для решения поставленных задач используются методы математического и функционального анализа (матричное дифференцирование, формула Тейлора, спектральный анализ симметричных матриц), методы теории вероятностей (анализ моментов гауссовских квадратичных форм, метод Монте-Карло), а также методы вычислительной линейной алгебры (произведения матрицы Гессе на вектор, дефлированная степенная итерация для построения собственных подпространств).

\textbf{Научная новизна.}

\begin{enumerate}
    \item Предложен унифицированный критерий $\Delta_p$ стабилизации поверхности функции потерь с произвольной функцией предпочтения точек в пространстве параметров, охватывающий одноточечный и изотропный среднеквадратичный случаи и формализующий выбор пробной области.
    \item Доказана верхняя оценка $\mathcal{O}(k^{-2})$ для среднеквадратичного критерия $\Delta_2^{(D)}$, ограниченного на $D$-мерное подпространство ведущих собственных векторов матрицы Гессе; в полученной оценке зависимость от размерности пространства параметров заменена на зависимость от размерности подпространства~$D$. Установлены замкнутое спектральное представление этого критерия и экстремальность подпространства главных направлений в семействе равноразмерных собственных подпространств.
    \item Разработаны масштабируемые алгоритмы оценивания критериев на основе произведений матрицы Гессе на вектор и замкнутого выражения через моменты гауссовских квадратичных форм, позволяющие вычислять критерии без явного построения матрицы Гессе.
    \item На основе доказанных оценок скорости убывания критериев получено количественное выражение для достаточного размера обучающей выборки, обеспечивающего заданный уровень стабилизации поверхности функции потерь.
\end{enumerate}

\textbf{Положения, выносимые на защиту.}

\begin{enumerate}
    \item Введенное семейство критериев $\Delta_p$ с функцией предпочтения точек в пространстве параметров охватывает как одноточечную, так и интегральные постановки задачи о стабилизации поверхности функции потерь и позволяет выделить пробную область, наиболее информативную с точки зрения локальной геометрии минимума.
    \item Среднеквадратичный критерий, ограниченный на $D$-мерное подпространство ведущих собственных векторов матрицы Гессе, сохраняет скорость убывания $\mathcal{O}(k^{-2})$ полнопространственного критерия, заменяя зависимость от размерности пространства параметров на зависимость от размерности подпространства, и допускает замкнутое спектральное выражение; подпространство главных направлений Гессе является экстремальным в семействе равноразмерных собственных подпространств.
    \item Предложенные алгоритмы оценивания на основе произведений матрицы Гессе на вектор и замкнутого выражения через моменты гауссовских квадратичных форм позволяют вычислять критерии стабилизации в моделях большой размерности без явного построения матрицы Гессе и согласуются с прямой оценкой методом Монте-Карло в пределах численной точности.
    \item Получаемые оценки скорости убывания критериев дают количественное выражение для достаточного размера обучающей выборки, обеспечивающего заданный уровень стабилизации поверхности функции потерь.
\end{enumerate}

\textbf{Теоретическая значимость} полученных результатов состоит в том, что предложенная схема дает строгую количественную связь между объемом обучающей выборки и локальной геометрией поверхности функции потерь в окрестности минимума, согласованную с эмпирически наблюдаемой анизотропией спектра матрицы Гессе нейронных сетей. Эта связь дополняет известные эмпирические законы масштабирования математически обоснованной оценкой скорости стабилизации поверхности функции потерь и количественным критерием достаточного размера выборки.

\textbf{Практическая значимость} настоящей работы заключается в возможности применения разработанных методов для обоснованного определения достаточного объема данных при планировании экспериментов и обучения параметрических моделей машинного обучения, а также для диагностики стабилизации модели на конкретной обучающей выборке. Подпространственный критерий $\Delta_2^{(D)}$ и сопровождающие алгоритмы оценивания позволяют применять полученные оценки в нейронных сетях современных размеров, где изотропные оценки расходятся из-за размерностного множителя $N^2$. Замкнутая форма оценки через моменты квадратичной формы гауссовского вектора экспериментально оказывается приблизительно в $1{,}8\cdot 10^4$ раз быстрее прямой Монте-Карло оценки без потери точности в локальном режиме применимости.

\textbf{Достоверность результатов} обеспечивается строгими математическими доказательствами полученных утверждений в рамках четко сформулированных предположений, а также экспериментальной проверкой теоретических оценок на полносвязной нейронной сети и трансформерной языковой модели с воспроизводимыми протоколами и публично доступной программной реализацией.

\textbf{Апробация результатов.} Основные результаты работы изложены в препринте~\cite{kiselev2025curvature} <<Curvature-Aligned Probing for Local Loss-Landscape Stabilization>>, поданном на конференцию NeurIPS~2026, и в публикациях~\cite{kiselev2024unraveling,kiselev2025ssdlb,kiselev2025ssdpdp,meshkov2024convnets}. Отдельные результаты докладывались на следующих научных конференциях:
\begin{enumerate}
    \item Открытая конференция ИСП РАН, Москва, 9--10~декабря 2025~г.
    \item 22-я Всероссийская конференция с международным участием <<Математические методы распознавания образов>> (ММРО"=2025), Муром, 22--26~сентября 2025~г.~\cite{kiselev2025mmro}.
    \item 67-я Всероссийская научная конференция МФТИ, Москва, 31 марта~--- 5~апреля 2025~г.~\cite{kiselev2025conf67mipt,meshkov2025conf67mipt}.
    \item Открытая конференция ИСП РАН, Москва, 11--12~декабря 2024~г.
    \item 66-я Всероссийская научная конференция МФТИ, Москва, 1--6~апреля 2024~г.~\cite{kiselev2024conf66mipt}.
\end{enumerate}

\textbf{Личный вклад автора.} Все приведенные в работе теоретические результаты и вычислительные эксперименты, кроме отдельно оговоренных случаев, получены автором лично при научном руководстве к.ф.-м.н.~А.\,В.~Грабового.

\textbf{Структура и объем работы.} Работа состоит из введения, четырех глав, заключения, списка использованных источников и приложения. В главе~\ref{sec1} проводится обзор существующих подходов к измерению стабилизации ландшафта функции потерь и формулируется research gap. В главе~\ref{sec2} вводится унифицированный критерий $p$-сходимости с функцией предпочтения точек в пространстве параметров. В главе~\ref{sec3} доказываются основные теоретические результаты "--- оценки скорости убывания критериев $\Delta_1$, $\Delta_2$ и $\Delta_2^{(D)}$ при росте объема обучающей выборки, формулируется утверждение о $\|\mathbf G\|_2$ для полносвязных нейронных сетей и описываются масштабируемые алгоритмы оценивания. В главе~\ref{sec4} проведенные теоретические оценки эмпирически проверяются на полносвязной нейронной сети и трансформерной языковой модели. % TODO: после финальной сборки указать общий объем в страницах,
% число рисунков, таблиц и количество наименований в списке литературы.
